\documentclass[twocolumn]{aastex631}
\begin{document}
\title{The reionization ionizing-photon-budget shortfall for star-forming galaxies is not robust to systematics at z\~{}7 (highC systematic corner)}
\author{NebulaMind Lab (autonomous pipeline)}
\affiliation{NebulaMind Lab, \url{https://lab.nebulamind.net}}
\begin{abstract}
Reionization ionizing-photon-budget reconciliation at z\~{}7: star-forming galaxies require f\_esc=0.184 (+0.173/-0.091) to close the budget (Madau-Dickinson SFRD, log xi\_ion=25.5+/-0.15, clumping C=2-5, JWST-SFRD tail), versus indirect-proxy-inferred f\_esc=0.062 (+0.108/-0.039) from LzLCS O32/beta calibrations. Median delta(required-inferred)=+0.107 dex-frac (16-84\%: -0.016 to +0.281); 81\% of the systematic MC shows a shortfall. Conclusion: the budget CLOSES within the systematic, and the sign holds under both O32 and beta calibrations. The result is bounded by the xi\_ion x clumping x proxy-calibration systematic, not by statistics. Generated autonomously from public data (jwst) via the NebulaMind Lab runner: a bounded, reproducible, descriptive study.
\end{abstract}
\section{Introduction}
Recent studies have highlighted a potential challenge in reconciling the observed ionizing photon production from star-forming galaxies with the requirements of cosmic reionization [Muñoz2024]. This discrepancy has sparked interest in revisiting the assumptions underlying these calculations and exploring ways to refine our understanding of the ionizing photon budget. Previous work has emphasized the importance of accurately accounting for factors such as the clumping factor, escape fraction, and ionizing efficiency when assessing the ionizing photon budget [Park2022, Davies2021].
\section{Data and method}
To address this issue, we employ a literature-anchored approach that relies on published values rather than new observational data. Specifically, we adopt the cosmic star formation rate density (SFRD) from Madau \& Dickinson (2014), as well as the ionizing efficiency (xi\_ion) and escape fraction (f\_esc) proxy calibrations derived from the LzLCS survey [Chisholm+22, Flury+22]. Our method involves calculating the reionization photon budget using these parameters and comparing it to the required budget based on indirect proxies. However, we acknowledge that this approach may oversimplify the complexity of key quantities like escape fraction and clumping factor due to unaccounted systematic uncertainties in proxy calibrations (e.g., LzLCS O32/beta).
\section{Result}
Our analysis suggests that star-forming galaxies at z\~{}7 require an escape fraction of f\_esc = 0.184 (+0.173/-0.091) to reconcile the ionizing photon budget, assuming a Madau-Dickinson SFRD, log xi\_ion = 25.5 +/- 0.15, and clumping factor C=2-5. This value is higher than the indirect-proxy-inferred escape fraction of f\_esc = 0.062 (+0.108/-0.039) derived from LzLCS O32/beta calibrations [Chisholm+22, Flury+22]. The median difference between the required and inferred escape fractions is +0.107 dex-frac (16-84\%: -0.016 to +0.281), with 81\% of systematic Monte Carlo simulations showing a shortfall. However, we caution that these results are subject to uncertainties in the adopted parameter values and their associated assumptions.
\begin{figure}\centering\includegraphics[width=\columnwidth]{result.png}\caption{The reionization ionizing-photon-budget shortfall for star-forming galaxies is not robust to systematics at z\~{}7 (highC systematic corner)}\end{figure}
\section{Caveats}
It is essential to acknowledge the limitations of our approach. Our analysis relies on a single selection of published values for key parameters, which may not fully capture the complexity and variability of these quantities in the real universe. Additionally, our method does not account for potential systematic uncertainties in the calibration of indirect proxies or the assumptions underlying the Madau-Dickinson SFRD model. Furthermore, our results are sensitive to the choice of clumping factor, which remains poorly constrained by current observations. These caveats highlight the need for further research and improved observational constraints to refine our understanding of the reionization photon budget. Future work should aim to incorporate a broader range of parameter values and their associated uncertainties, especially for the clumping factor, to provide a more robust assessment of the photon budget crisis.
\section*{References}
\footnotesize
[Muoz2024] Muñoz et al. (2024). Reionization after JWST: a photon budget crisis?. bibcode 2024MNRAS.535L..37M \par [Park2022] Park et al. (2022). Calibrating excursion set reionization models to approximately conserve ionizing photons. bibcode 2022MNRAS.517..192P \par [Duncan2015] Duncan et al. (2015). Powering reionization: assessing the galaxy ionizing photon budget at z < 10. bibcode 2015MNRAS.451.2030D \par [Davies2021] Davies et al. (2021). The Predicament of Absorption-dominated Reionization: Increased Demands on Ionizing Sources. bibcode 2021ApJ...918L..35D \par [Madau2017] Madau (2017). Cosmic Reionization after Planck and before JWST: An Analytic Approach. bibcode 2017ApJ...851...50M

\end{document}
